% =========================================================
%  sesi/01/sesi.tex
%  Sesi 1 — Lingkungan kerja dan jendela pertama
% =========================================================

\sesi{01}{Lingkungan Kerja dan Jendela Pertama}{1.1 (persiapan)}{100 menit}

Anda akan membangun \NamaPermainan: sebuah gudang tua berisi artefak. Pemain
masuk, mengumpulkan artefak, menghindari peti dan penjaga, lalu keluar lewat
pintu. Permainan ini tidak dibuat sekaligus. Setiap sesi menambahkan satu
kemampuan, dan pada sesi terakhir seluruh kemampuan itu bekerja bersama
dalam satu program.

Sesi pertama belum membuat permainan. Sesi ini menyiapkan tempat kerja dan
membuktikan bahwa perkakas Anda benar-benar berjalan: sebuah jendela kosong
berjudul \NamaPermainan.

% ---------------------------------------------------------
\bagian{Tujuan Sesi}

\begin{tujuanbox}
Setelah menuntaskan sesi ini, Anda dapat menunjukkan:
\begin{enumerate}[leftmargin=1.5em, itemsep=0.2em]
  \item jendela berukuran 960 $\times$ 600 piksel dengan judul
        \NamaPermainan\ terbuka di layar;
  \item lingkungan kerja yang memuat Python 3.8 atau lebih baru dan Pygame 2;
  \item folder proyek dengan paket \file{gudang_artefak} dan berkas
        \file{main.py} yang dapat dijalankan;
  \item jendela yang dapat ditutup dengan tombol silang tanpa meninggalkan
        proses yang menggantung;
  \item penjelasan singkat mengapa program berhenti bersih setelah
        \code{pygame.quit()}.
\end{enumerate}
\end{tujuanbox}

% ---------------------------------------------------------
\bagian{Prasyarat}

\begin{prasyaratbox}
\begin{itemize}[leftmargin=1.4em, itemsep=0.25em]
  \item Anda pernah menulis program Python sederhana: variabel, fungsi, dan
        kelas.
  \item Anda dapat membuka terminal dan berpindah folder dengan \code{cd}.
  \item Tidak ada checkpoint sebelumnya. Bab I dan Bab II buku ajar cukup
        dibaca sebagai orientasi: keduanya menjelaskan waktu, ruang, dan
        pengendalian objek yang akan Anda pakai nanti, tetapi tidak perlu
        dihafal sekarang.
\end{itemize}
\end{prasyaratbox}

% ---------------------------------------------------------
\bagian{Peta Kemampuan}

\petasesi
  {Belum ada program.}
  {Jendela berjudul, paket \file{gudang_artefak}, dan titik masuk
   \file{main.py}.}
  {Siklus permainan yang berjalan terus-menerus dan pengendali waktu pada
   sesi 2.}

% ---------------------------------------------------------
\bagian{Konsep yang Dipakai}

\subsection*{Pygame adalah pustaka, bukan bahasa}

Python yang Anda tulis dijalankan oleh penafsir. Pygame bukan bahasa lain dan
bukan program berdiri sendiri; ia kumpulan fungsi yang dipanggil dari program
Python untuk membuka jendela, menggambar, memuat bunyi, dan membaca masukan.
Di bawah Pygame bekerja SDL yang berbicara dengan perangkat keras.

\begin{center}
% figures/lapisan-pygame.tikz
% Lapisan perkakas: program, Pygame, SDL, perangkat.
\begin{tikzpicture}[
  lapis/.style={rectangle, rounded corners=2pt, draw=ubNavy, thick,
    fill=ubSoft, align=center, text width=7.4cm, inner sep=4pt,
    minimum height=0.95cm, font=\small},
  sdl/.style={lapis, draw=ubAccent, fill=white},
  os/.style={lapis, draw=ubLine, fill=ubGray},
  panah/.style={-{Stealth[length=2.2mm]}, thick, draw=ubAccent}
]
  \node[lapis] (program) {\textbf{Program Anda}\\ \footnotesize kode pada paket gudang\_artefak};
  \node[sdl, below=0.55cm of program] (pygame) {\textbf{Pygame}\\ \footnotesize jendela, gambar, huruf, bunyi, peristiwa};
  \node[sdl, below=0.55cm of pygame] (sdl) {\textbf{SDL}\\ \footnotesize penerjemah ke perangkat keras};
  \node[os, below=0.55cm of sdl] (os) {\textbf{Sistem operasi dan perangkat}\\ \footnotesize layar, papan tombol, tetikus, pengeras suara};
  \draw[panah] (program) -- (pygame);
  \draw[panah] (pygame) -- (sdl);
  \draw[panah] (sdl) -- (os);
\end{tikzpicture}

\end{center}

Karena berlapis, ada dua akibat yang akan Anda temui. Pertama, pesan kesalahan
kadang menyebut SDL, bukan Pygame, karena masalahnya berada di lapisan yang
lebih dalam. Kedua, Pygame yang terpasang pada penafsir yang berbeda tidak
akan ditemukan oleh program Anda, walaupun pemasangan sebelumnya berhasil.

\subsection*{Kerangka program dalam lima langkah}

Setiap program Pygame, sederhana maupun rumit, mengikuti kerangka berikut.
Kelima langkah ini akan muncul kembali, dengan tambahan, pada setiap sesi.

\begin{enumerate}[leftmargin=1.6em, itemsep=0.25em]
  \item Memulai Pygame dengan \code{pygame.init()}.
  \item Membuat permukaan tampilan dengan \code{pygame.display.set_mode()},
        yang membuka jendela sekaligus mengembalikan tempat menggambar.
  \item Menggambar isi jendela: mengisi latar, lalu menempelkan bentuk atau
        teks.
  \item Menayangkan hasil gambar dengan \code{pygame.display.flip()}.
  \item Menunggu sampai program diakhiri, lalu membersihkan dengan
        \code{pygame.quit()}.
\end{enumerate}

\subsection*{Menggambar dan menayangkan adalah dua pekerjaan}

Apa pun yang digambar di Pygame --- jendela, bentuk, atau teks --- adalah
\istilah{permukaan} (\emph{surface}). Menggambar berarti menempelkan satu
permukaan ke permukaan lain. Hasilnya masih berada di memori dan belum
terlihat pemain sampai \code{flip()} dijalankan. Bila Anda mengubah warna
latar tetapi layar tidak berubah, periksa apakah \code{flip()} masih
dipanggil.

Sesi ini masih memakai perulangan peristiwa: program menunggu satu peristiwa,
lalu bereaksi. Perulangan itu belum siklus permainan, karena tidak ada
keadaan yang diperbarui dan tidak ada gambar yang disusun ulang.

% ---------------------------------------------------------
\bagian{Tutorial}

Kerjakan langkah berikut berurutan. Setiap langkah adalah satu tindakan yang
selesai; jangan melompati langkah hanya karena tampaknya sudah jelas.

\begin{langkah}
  \item \textbf{Periksa penafsir Python.} Buka terminal dan jalankan
        \code{python --version}. Bila perintah itu tidak dikenali, coba
        \code{py --version}. Yang dibutuhkan adalah Python 3.8 atau lebih
        baru; tampilan versi 3.11 berarti siap.

  \item \textbf{Pasang Pygame.} Jalankan
        \code{python -m pip install pygame}. Bentuk \code{python -m pip}
        memastikan paket dipasang pada penafsir yang sama dengan yang akan
        menjalankan program. Bila muncul pesan izin ditolak, ulangi dengan
        \code{python -m pip install --user pygame}.

  \item \textbf{Buat folder proyek.} Buat folder bernama
        \file{gudang-artefak}, lalu di dalamnya buat folder paket
        \file{gudang_artefak}. Nama paket memakai huruf kecil dan garis
        bawah, tanpa spasi, agar aman pada semua sistem operasi.

  \item \textbf{Buat penanda paket.} Di dalam \file{gudang_artefak}, buat
        berkas kosong bernama \file{__init__.py}. Berkas ini menandai bahwa
        folder itu adalah paket Python; isinya boleh hanya satu baris
        keterangan.

  \item \textbf{Tulis pengaturan.} Buat \file{gudang_artefak/settings.py}
        berisi judul, ukuran jendela, dan warna latar. Semua angka yang
        dipakai bersama dikumpulkan di sini supaya mudah dicari ketika
        terjadi masalah.

  \item \textbf{Tulis kelas permainan.} Buat
        \file{gudang_artefak/game.py}. Kelas \code{Game} membuka jendela di
        \code{__init__} dan menunggu peristiwa di \code{run()}.

  \item \textbf{Tulis titik masuk.} Buat \file{main.py} yang membaca pilihan
        baris perintah, membuat \code{Game}, lalu memanggil \code{run()}.
        Pilihan \code{--auto-close} disediakan untuk pengujian: program
        mengirim sendiri peristiwa penutup sehingga berhenti tanpa disentuh.

  \item \textbf{Jalankan program.} Dari folder \file{gudang-artefak}, ketik
        \code{python main.py}. Jendela biru tua berjudul \NamaPermainan\
        harus terbuka.
\end{langkah}

\begin{ujibox}
Setelah Langkah 8, periksa empat hal berikut.
\begin{itemize}[leftmargin=1.4em, itemsep=0.2em]
  \item Jendela terbuka dengan ukuran 960 $\times$ 600 piksel. Bila berbeda,
        periksa nilai bawaan pada \code{parse_arguments}.
  \item Judul pada bilah jendela berbunyi \NamaPermainan. Bila masih bawaan
        Pygame, periksa \code{pygame.display.set_caption}.
  \item Latar jendela berwarna biru tua, bukan hitam. Bila hitam, periksa
        urutan \code{fill} dan \code{flip}.
  \item Menekan tombol silang menutup jendela, dan terminal menampilkan
        \code{Program selesai.} Bila terminal tidak kembali, periksa
        pemanggilan \code{pygame.quit()}.
\end{itemize}
\end{ujibox}

% ---------------------------------------------------------
\bagian{Kode Sesi Ini}

Ketiga berkas berikut adalah seluruh isi checkpoint sesi 1. Anda tidak perlu
mengetik ulang kode di bawah dari layar; ketiklah sendiri di editor agar
setiap baris terbaca, lalu bandingkan hasilnya dengan daftar ini.

\tampilkode[language=Python]{code/sesi-01/gudang_artefak/settings.py}
  {settings.py: pengaturan yang dipakai bersama}

\tampilkode[language=Python]{code/sesi-01/gudang_artefak/game.py}
  {game.py: membuka jendela dan menunggu peristiwa}

\tampilkode[language=Python]{code/sesi-01/main.py}
  {main.py: membaca pilihan baris perintah dan menjalankan permainan}

\checkpoint{01}{Program pada folder itu dapat dijalankan langsung dengan
\code{python main.py} dan berhenti dengan rapi ketika jendela ditutup.}

% ---------------------------------------------------------
\bagian{Kesalahan Umum}

\begin{galatbox}
\begin{itemize}[leftmargin=1.4em, itemsep=0.35em]
  \item \textbf{\code{ModuleNotFoundError: No module named 'pygame'}.}
        Gejala: program berhenti sebelum jendela terbuka. Periksa baris
        \emph{Penafsir Python} yang dipakai; bila ada dua versi Python,
        pasang ulang dengan \code{python -m pip install pygame} dari terminal
        yang sama dengan yang menjalankan program.
  \item \textbf{Jendela terbuka lalu langsung tertutup.}
        Gejala: tidak ada kesalahan, tetapi jendela hanya berkedip. Periksa
        apakah program dijalankan dengan \code{--auto-close}, atau apakah
        perulangan peristiwa berakhir karena variabel \code{running} sudah
        bernilai salah.
  \item \textbf{\code{IndentationError: unexpected indent}.}
        Gejala: program tidak pernah sampai membuka jendela. Periksa
        indentasi di sekitar baris yang ditunjuk; jangan mencampur tab dan
        spasi. Aturlah editor mengubah tab menjadi empat spasi.
  \item \textbf{\code{FileNotFoundError: [Errno 2]}.}
        Gejala: Python menyebut berkas tidak ditemukan padahal berkasnya ada.
        Periksa folder kerja terminal dengan \code{dir} atau \code{ls};
        jalankan \code{cd} ke folder yang benar.
  \item \textbf{Layar hitam setelah \code{set_mode}.}
        Gejala: jendela terbuka tetapi isinya hitam. Periksa urutan
        \code{fill}, lalu \code{flip}; menayangkan sebelum menggambar
        menghasilkan layar kosong.
\end{itemize}
\end{galatbox}

% ---------------------------------------------------------
\bagian{Lembar Kerja}

Isi tabel berikut pada catatan praktikum Anda. Kolom terakhir diisi setelah
Anda mencoba mengubah satu nilai dan menjalankan ulang program.

{\small
\renewcommand{\arraystretch}{1.4}
\begin{tabularx}{\linewidth}{@{}C{0.8cm} L{5.2cm} Y@{}}
\toprule
\textbf{No} & \textbf{Hasil yang dicatat} & \textbf{Keterangan Anda} \\
\midrule
1 & Versi Python dan versi Pygame & \\[0.55cm]
2 & Nama berkas penafsir Python yang dipakai & \\[0.55cm]
3 & Ukuran jendela dan judul yang tampil & \\[0.55cm]
4 & Nilai \code{BACKGROUND} sebelum dan sesudah diubah & \\[0.55cm]
5 & Satu kendala yang Anda hadapi dan cara menyelesaikannya & \\[0.9cm]
\bottomrule
\end{tabularx}
}

% ---------------------------------------------------------
\bagian{Checklist Sesi}

\begin{checklistbox}
\begin{checklist}
  \item Saya dapat menunjukkan jendela \NamaPermainan\ yang terbuka dengan
        ukuran yang saya tentukan sendiri.
  \item Saya dapat menjelaskan perbedaan peran Python dan Pygame.
  \item Saya dapat menyebutkan lima langkah kerangka program Pygame.
  \item Saya dapat menjelaskan mengapa gambar baru terlihat setelah
        \code{flip()}.
  \item Saya dapat menutup program dengan tombol silang dan memastikan
        terminal kembali dengan rapi.
  \item Saya dapat membaca pesan kesalahan dan menentukan baris penyebabnya.
\end{checklist}
\end{checklistbox}

% ---------------------------------------------------------
\bagian{Pengait ke Sesi Berikutnya}

Jendela pada sesi ini hanya menunggu peristiwa; isinya tidak pernah berubah.
Pada sesi 2 Anda menambahkan siklus permainan: membaca masukan, memperbarui
keadaan, menggambar ulang, lalu menayangkannya berulang kali. Di sanalah
\istilah{Clock} mulai mengatur laju bingkai dan angka FPS muncul pada judul
jendela.
